\documentclass[10pt,a4paper]{amsart}
\usepackage[margin=19mm]{geometry}
\usepackage{amsmath,amssymb,mathtools}
\usepackage[scr=rsfs,cal=euler]{mathalfa}
\usepackage{xfrac}
\usepackage[unicode,hidelinks,bookmarks=false]{hyperref}
\newcommand{\ie}{i.e.\ }
\newcommand{\cf}{cf.\ }
\newcommand{\resp}{respectively\ }
\newcommand{\Subsection}{\subsection}
\providecommand{\nobreakdash}{\nobreak\hbox{-}}
\setlength{\emergencystretch}{2em}
\multlinegap0pt
\usepackage{amssymb}
\usepackage{dsfont}
\usepackage{mathtools}
\usepackage{xcolor}
\usepackage{float}
\usepackage{tikz}
\newtheorem{lemma}{Lemma}
\def\R{\mathbb{R}}
\def\d{\mathrm{d}}
\def\exp{\mathrm{exp}}
\def\no{\nonumber}
\title{Global renormalized solutions to Boltzmann systems modeling mixture gases of monatomic and polyatomic species}
\author{Yi-Long Luo, Jing-Xin Nie${}^*$}
\date{Source: arXiv:2601.07480v1 (CC BY 4.0)}
\begin{document}
\maketitle

\noindent\emph{Source and licence.} Global renormalized solutions to Boltzmann systems modeling mixture gases of monatomic and polyatomic species. Version arXiv:2601.07480v1 (CC BY 4.0), distributed by arXiv under the Creative Commons Attribution 4.0 International licence, http://creativecommons.org/licenses/by/4.0/, as recorded in the arXiv licence metadata for this exact version and shown on the abstract page https://arxiv.org/abs/2601.07480v1. Full licence notice: \texttt{licenses/arxiv-2601.07480-CC-BY-4.0.txt}.\par\smallskip

\noindent\textit{Editorial scope.} Counted unit: the article's lemma L (source0.tex lines 2853--2865). The article's complete original proof of it is delivered with it (source0.tex lines 2868--3009, one \texttt{proof} environment of 142 lines). No mathematics is written, repaired or re-derived: the statement and the proof are the article's own lines, in the article's own notation, and no retained line is substituted or re-flowed.  Retained uncounted as the source-local context the proof needs: lines 415--419; lines 428--432; lines 441--445; lines 456--461; lines 505--513; lines 593--595; lines 606--613; lines 648--653; lines 1193--1195; lines 1202--1215; lines 1298--1311; lines 1414--1435; lines 1496--1527; lines 1581--1596; lines 1704--1722; lines 1724--1728; lines 1989--2002; lines 2130--2213; lines 2325--2331; lines 2526--2532; lines 2535--2541; lines 2545--2552; lines 2555--2561; lines 2597--2613.  Nothing was omitted from any retained span, so the package carries no omission record.  No retained line contains a citation, so the wrapper carries no bibliography and no bibliography entry.  No retained line contains a figure environment, an image-inclusion command or drawing code, so no image is delivered and the package declares no asset dependency.  Editorial formatting only, in addition: an AMS \texttt{amsart} preamble that restates the article's own declarations and macros used by the retained lines, namely the environments \texttt{lemma} on separate counters, together with the standard packages those lines need.  The retained environments print in this document as Lemma 1 in order of appearance, whereas the article prints the same items with the numbering of its own full document.  The complete native source of the article as distributed in the arXiv e-print for this exact version is bundled unchanged as \texttt{sources/192504/source0.tex}.
    \begin{equation}\label{B0-ab}
        \begin{aligned}
            B_{0\alpha \beta } (\xi-\xi_*,\omega )=\tilde{\sigma}_{\alpha \beta }|g|,\quad \tilde{\sigma}_{\alpha \beta }(|g|,\omega)=\sigma_{\alpha \beta }(|g|,|\cos \theta|) \,;
        \end{aligned}
    \end{equation}

    \begin{equation}\label{B1-ab}{\small
        \begin{aligned}
            B_{1\alpha \beta }(\xi-\xi_*,I_*,\mathfrak{R},\omega )  = \tfrac{\tilde{\sigma}_{\alpha \beta }|g|E_{\alpha \beta }}{\mathfrak{R}^{1/2} (1-\mathfrak{R})^{\delta(\beta)/2-1}},\  \tilde{\sigma}_{\alpha \beta } (|g|,I_* ,\mathfrak{R},\omega )=\sigma_{\alpha \beta }(|g|,|\cos \theta|,I_*,I'_*) \,;
        \end{aligned}}
    \end{equation}

    \begin{equation}\label{B1-ba}
        \begin{aligned}
            B_{1 \beta \alpha }(\xi-\xi_*,I,\mathfrak{R},\omega )  = \tfrac{\tilde{\sigma}_{\alpha \beta }|g|E_{\alpha \beta }}{\mathfrak{R}^{1/2} (1-\mathfrak{R})^{\delta(\alpha)/2-1}},\  \tilde{\sigma}_{\alpha \beta } (|g|,I ,\mathfrak{R},\omega )=\sigma_{\alpha \beta }(|g|,|\cos \theta|,I,I') \,;
        \end{aligned}
    \end{equation}

    \begin{equation}\label{B2-ab}
        \begin{aligned}
            B_{2\alpha \beta }(\xi-\xi_*,I+I_*,\mathfrak{R},\mathfrak{r},\omega)&= \tfrac{\tilde{\sigma}_{\alpha \beta } |g|E^2_{\alpha \beta }}{\mathfrak{r}^{\delta(\alpha)/2-1} (1-\mathfrak{r})^{\delta (\beta )/2-1 }(1-\mathfrak{R} ) ^{(\delta (\alpha) +\delta (\beta ))/2-2} \mathfrak{R}^{1/2} } \,,\\
            \tilde{\sigma}_{\alpha \beta } (|g|,I+I_*,\mathfrak{R},\mathfrak{r},\omega) &= \sigma_{\alpha \beta } (|g|,|\cos \theta|,I+I_*,I'+I'_*) 
        \end{aligned}
    \end{equation}

\begin{equation}\label{A}{\small
    \begin{aligned}
		A_{0 \alpha \beta } (\xi-\xi_*)=&\int_{\mathbb{S}^2 } B_{0 \alpha \beta }(\xi-\xi_*,\omega ) d\omega, \\
		A_{1\alpha \beta }(\xi-\xi_*,I_*)= &\iint_{[0,1] \times \mathbb{S}^2 } B_{1 \alpha \beta } (\xi-\xi_*,I_*,\mathfrak{R}, \omega )(1-\mathfrak{R})^{\delta(\beta )/2-1}\mathfrak{R}^{1/2} d\mathfrak{R}d\omega, \\
		A_{1 \beta \alpha }(\xi-\xi_*,I)=&\iint_{[0,1] \times \mathbb{S}^2 } B_{1 \beta \alpha }(\xi-\xi_*,I,\mathfrak{R},\omega) (1-\mathfrak{R})^{\delta(\alpha )/2-1}\mathfrak{R}^{1/2} d\mathfrak{R}d\omega, \\
		A_{2\alpha \beta } (\xi-\xi_*,I+I_*)=&\iiint_{[0,1]^2 \times \mathbb{S}^2 } B_{2\alpha \beta }(\xi-\xi_*,I+I_*,\mathfrak{R},\mathfrak{r},\omega ) \\
     & \times \mathfrak{r}^{\delta(\alpha )/2-1} (1-\mathfrak{r})^{\delta(\beta )/2-1} (1-\mathfrak{R})^{(\delta(\alpha)+\delta (\beta ))/2-1} \mathfrak{R}^{1/2} d\mathfrak{r}d\mathfrak{R}d\omega .
    \end{aligned}}
\end{equation}

        \begin{equation}\label{fs0}
            \iint_{\R^3 \times \R^3} f_{\alpha,0} (1+|x|^2 +|\xi|^2 +|\log f_{\alpha,0}|)dxd\xi <\infty \,.
        \end{equation}

\begin{equation}\label{collision}
    \begin{aligned}
        & ( 1 + | \xi |^2 )^{-1} \int_{ | z - \xi | \leq R } A_{ 0 \alpha \beta } (z) d z \rightarrow 0, \quad | \xi | \rightarrow \infty, \\
        & ( 1 + | \xi |^2 +  \eta )^{-1} \int_{ | z -\xi | \leq R } A_{ 1 \alpha \beta } ( z, \eta ) d z \rightarrow 0,  \quad | \xi |,\eta  \rightarrow \infty, \\
        & ( 1 + | \xi |^2 )^{-1} \int_{ | z - \xi | \leq R } \int_{ 0 < \eta < R } A_{ 1 \beta \alpha } ( z, \eta ) d z d \eta \rightarrow 0, \quad | \xi | \rightarrow \infty , \\
        & ( 1 + | \xi |^2 +  I )^{-1} \int_{ | z - \xi | \leq R } \int_{ 0 < \eta - I < R } A_{ 2 \alpha \beta } ( z , \eta ) d z d \eta \rightarrow 0, \quad | \xi |, I \rightarrow \infty .
    \end{aligned}
\end{equation}

\begin{equation}\label{Hf}
	\begin{aligned}
		H(f)(t)& = (f,\log (\varphi^{-1} f) ) \\
		&=\sum \limits_{\alpha=1}^{s_0} \iint_{\R^3 \times \R^3 } f_\alpha \log f_\alpha  dxd\xi + \sum \limits_{\alpha = s_0+1}^{s}\iiint_{\R^3 \times \R^3 \times \R_+ } f_\alpha \log (I^{1-\delta(\alpha)/2} f_\alpha ) dxd\xi dI.
	\end{aligned}
\end{equation}

\begin{equation} \label{approximation}
		\tfrac{\partial f^n}{\partial t} +\xi \cdot \nabla_x f^n =\tilde{Q}_n(f^n, f^n) \,, \ f^n|_{t=0}=f^n_{0} \,,
\end{equation} 

\begin{lemma}\label{Lemma 4.8}
	Assume that $f_{\alpha,0 } \in L^1_2 ( \R^3 \times \R^3)_+ $ for $\alpha \in \{1,\cdots ,s_0\}$ satisfy \eqref{fs0}, hence,
	\begin{equation*}
		\iint_{\R^3 \times \R^3 } f_{\alpha,0 } (x,\xi ) \cdot (1+|x|^2 +|\xi|^2 + |\log  f_{\alpha,0 } |)  dxd\xi <+ \infty .
	\end{equation*}
	Then there exists a sequence $\{f^n_{\alpha,0 }\}_{n\geq 1} \subset \mathscr{S} (\R^3 \times \R^3 )$ (the Schwartz space over $\R^3 \times \R^3$) and a constant $C > 0$ such that 
\begin{equation}\label{Ap-fs0}
  \begin{aligned}
    & \lim \limits_{n\rightarrow \infty}  \iint_{\R^3 \times \R^3 } (1+|x|^2 +|\xi|^2) |f^n_{\alpha,0 } - f_{\alpha,0 }| dxd\xi=0 \,, \ f^n_{\alpha,0 } \geq \tfrac{1}{n} \exp (- \tfrac{|x|^2+ |\xi|^2 }{2}) \,, \\
    & \iint_{\R^3 \times \R^3} f^n_{\alpha,0 }  (1+|x|^2 +|\xi|^2 + |\log f^n_{\alpha,0 }|)dxd\xi  \leq C \,, \\
    & \limsup \limits_{n\rightarrow \infty} \iint_{\R^3 \times \R^3 } f^n_{\alpha,0 } \log f^n_{\alpha,0 }  dxd\xi  \leq \iint_{\R^3 \times \R^3 } f_{\alpha,0 } \log  f_{\alpha,0 } dxd\xi \,.
  \end{aligned}
\end{equation}
\end{lemma}

\begin{lemma}\label{Lemma 4.9}
	Assume that $f_{\alpha,0 } \in L^1_3( \R^3 \times \R^3 \times \R_+)_+ $ for $\alpha \in \{s_0+1,\cdots ,s\}$ satisfy
	\begin{equation}\label{eq 4.51}
		\iiint_{\R^3 \times \R^3 \times \R_+ } f_{\alpha,0 } (1+|x|^2 +|\xi|^2+I + |\log (I ^{1-\frac{\delta(\alpha )}{2} } f_{\alpha,0 } )|)  dxd\xi dI<+ \infty .
	\end{equation}
	Then there exists a sequence $\{f^n_{\alpha,0 }\}_{n\geq 1} \subset \mathscr{S} (\R^3 \times \R^3 \times \R_+ ) $ (the Schwartz space over $\R^3 \times \R^3 \times \R_+$) and a constant $C > 0$ such that 
	\begin{enumerate}
		\item {$\lim \limits_{n\rightarrow \infty}  \iiint_{\R^3 \times \R^3 \times \R_+ } (1+|x|^2 +|\xi|^2+I) |f^n_{\alpha,0 } - f_{\alpha,0 }| dxd\xi dI=0 $;}\label{1 of Lemma 4.9}
		\item {$f^n_{\alpha,0 } \geq \frac{1}{n} \exp (- \frac{|x|^2+ |\xi|^2 + I }{2})$; }\label{2 of Lemma 4.9}
		\item {$\iiint_{\R^3 \times \R^3 \times \R_+ } f^n_{\alpha,0 } (1+|x|^2 +|\xi|^2 + I + |\log (I ^{1-\frac{\delta(\alpha )}{2} }f^n_{\alpha,0 })|)dxd\xi dI \leq C $;}\label{3 of Lemma 4.9}
		\item {\small $\limsup \limits_{n\rightarrow \infty} \iiint_{\R^3 \times \R^3 \times \R_+ } f^n_{\alpha,0 } \log (I ^{1-\frac{\delta(\alpha )}{2} } f^n_{\alpha,0 } ) dxd\xi dI
			\leq \iiint_{\R^3 \times \R^3 \times \R_+ } f_{\alpha,0 } \log (I ^{1-\frac{\delta(\alpha )}{2} } f_{\alpha,0 } ) dxd\xi dI$.}\label{4 of Lemma 4.9}
	\end{enumerate}
\end{lemma}

	\begin{lemma}\label{Lemma 4.10}
	Assume collision kernel $A_{2\alpha \beta }$ satisfies the assumption in \eqref{collision}, i.e.,
	\begin{equation*}
		(1+|\xi|^2 + I )^{-1} \int_{|z-\xi |\leq R } \int_{0<\eta -I < R } A_{2\alpha \beta } (z,\eta ) dzd\eta \rightarrow 0, \quad |\xi|, I \rightarrow \infty , \ \forall R <+\infty .
	\end{equation*}	
	Then there exists $B_{2\alpha \beta ,n }\geq 0 $ satisfying
	\begin{enumerate}
		\item {For any fixed $(\mathfrak{R},\mathfrak{r},\omega)\in [0,1]^2 \times \mathbb{S}^2, B_{2\alpha \beta , n } (z,\eta , \mathfrak{R},\mathfrak{r},\omega ) \in C^\infty (\R^3 \times \R_+), supp B_{2\alpha \beta , n } \subset \{(z,\eta , \mathfrak{R},\mathfrak{r},\omega ): \frac{1}{n} \leq |z|,\eta \leq n \}$ and if $ z \cdot \omega \leq \frac{1}{n}, B_{2\alpha \beta , n } =0$;	}\label{1 of Lemma 4.10}
		\item  {For almost all $(z,\eta , \mathfrak{R},\mathfrak{r},\omega )\in \R^3 \times \R_+ \times [0,1]^2 \times \mathbb{S}^2$, 
  $ \lim \limits_{n\rightarrow \infty } B_{2\alpha \beta ,n } =B_{2\alpha \beta } $, where $B_{2\alpha \beta }$ is defined in \eqref{B2-ab};
  } \label{2 of Lemma 4.10}
		\item {For any $R < + \infty$,
		\begin{equation*}
		\sup \limits_{n \geq 1} (1+|\xi|^2 + I )^{-1} \int_{|z-\xi |\leq R } \int_{0<\eta -I < R } A_{2\alpha \beta ,n } (z,\eta ) dzd\eta \rightarrow 0, \ |\xi|, I \rightarrow \infty ,
		\end{equation*}	
  where 
  \begin{equation*}
    A_{2\alpha \beta ,n }=\int_{[0,1]^2 \times \mathbb{S}^2} B_{2\alpha \beta ,n } \mathfrak{r}^{\delta(\alpha)/2-1 } (1-\mathfrak{r})^{\delta(\beta )/2-1} (1-\mathfrak{R})^{(\delta(\alpha)+(\delta(\beta ))/2-1} \mathfrak{R}^{1/2}  d\mathfrak{R}d\mathfrak{r}d\omega.  
  \end{equation*} 
}\label{3 of Lemma 4.10}
	\end{enumerate}
\end{lemma}

\begin{lemma}\label{Lemma 4.11}
	Assume that collision kernels $A_{1\alpha \beta }$ and $A_{1\beta \alpha }$ satisfy the assumptions \eqref{collision}, i.e.,
	\begin{equation*}{\small
		\begin{aligned}
			&(1+|\xi|^2 + \eta )^{-1} \int_{|z-\xi |\leq R } A_{1\alpha \beta } (z,\eta ) dz,  \quad |\xi|,\eta  \rightarrow \infty, \ \forall R <+\infty, \\
			&(1+|\xi|^2 )^{-1}\int_{|z-\xi |\leq R }  \int_{0<\eta < R } A_{1\beta \alpha } (z,\eta ) dzd\eta \rightarrow 0, \quad |\xi| \rightarrow \infty , \ \forall R <+\infty.
		\end{aligned}}
	\end{equation*}	
	Then there exist $B_{1\alpha \beta ,n }(z,\eta , \mathfrak{R},\omega ) ,B_{1 \beta \alpha ,n } (z,\eta , \mathfrak{R},\omega ) \geq 0 $ satisfying 
	\begin{enumerate}
		\item {For any $(\mathfrak{R},\omega)\in [0,1] \times \mathbb{S}^2,\ B_{1\alpha \beta , n } , B_{1 \beta \alpha ,n }  \in C^\infty (\R^3 \times \R_+),\ supp B_{1\alpha \beta , n },supp B_{1 \beta \alpha ,n } \subset \{(z,\eta , \mathfrak{R},\omega ): \frac{1}{n} \leq |z|,\eta \leq n \}$ and if $ (\xi -\xi_* ) \cdot \omega \leq \frac{1}{n} , B_{1\alpha \beta , n } =B_{1 \beta \alpha ,n }=0$;}\label{1 of Lemma 4.11}
		\item  {For almost all $(z,\eta , \mathfrak{R},\omega )\in \R^3 \times \R_+ \times [0,1] \times \mathbb{S}^2$,
		\begin{equation*}
  \lim \limits_{n\rightarrow \infty } B_{1\alpha \beta ,n } =B_{1\alpha \beta },\quad
 \lim \limits_{n\rightarrow \infty } B_{1\beta \alpha ,n } =B_{1\beta \alpha },
\end{equation*} 
where $B_{1 \alpha \beta}$ and $B_{1 \beta \alpha}$ are respectively defined in \eqref{B1-ab} and \eqref{B1-ba};} \label{2 of Lemma 4.11}
		\item {For all $R < + \infty$,
			\begin{equation*}{\small
				\begin{aligned}
					&\sup \limits_{n \geq 1} (1+|\xi|^2 + \eta )^{-1} \int_{|z-\xi |\leq R } A_{1\alpha \beta ,n } (z,\eta ) dz \rightarrow 0, \quad |\xi|, \eta \rightarrow \infty , \\
					&\sup \limits_{n \geq 1} (1+|\xi|^2 )^{-1}\int_{|z-\xi |\leq R }  \int_{0<\eta < R } A_{1\beta \alpha ,n} (z,\eta ) dzd\eta \rightarrow 0, \quad |\xi| \rightarrow \infty \,,
				\end{aligned}}
		\end{equation*}
   where 
{\small
    \begin{align*}
   A_{1\alpha \beta ,n } =\iint_{[0,1]\times \mathbb{S}^2} B_{1\alpha \beta ,n } (1-\mathfrak{R})^{\delta(\beta )/2-1}\mathfrak{R}^{1/2}  d\mathfrak{R}d\omega, A_{1\beta \alpha ,n } =\iint_{[0,1]\times \mathbb{S}^2} B_{1\beta \alpha ,n }   (1-\mathfrak{R})^{\delta(\alpha )/2-1}\mathfrak{R}^{1/2}d\mathfrak{R}d\omega.
   \end{align*}}
  }\label{3 of Lemma 4.11}
	\end{enumerate}
\end{lemma}

	\begin{lemma}\label{Lemma 4.12}
	Assume that collision kernel $A_{0\alpha \beta }$ satisfies the assumption \eqref{A}, hence,
	\begin{equation*}
		(1+|\xi|^2)^{-1} \int_{|z-\xi |\leq R } A_{0\alpha \beta } (z)dz, \quad |\xi|\rightarrow \infty, \ \forall R <+\infty .
	\end{equation*}	
	Then there exists $B_{0\alpha \beta ,n }\geq 0 $ satisfying
	\begin{enumerate}
		\item {For any $\omega\in \mathbb{S}^2, B_{0\alpha \beta ,n } (z,\omega ) \in C^\infty (\R^3 ),supp B_{0\alpha \beta , n} \subset \{(z,\omega ): \frac{1}{n} \leq |z| \leq n \}$ and if $ (\xi -\xi_* ) \cdot \omega \leq \frac{1}{n} , B_{0\alpha \beta ,n } =0$;}\label{1 of Lemma 4.12}
		\item  {For almost all $(z,\omega )\in \R^3 \times \mathbb{S}^2 , \lim \limits_{n\rightarrow \infty } B_{0\alpha \beta ,n } =B_{0\alpha \beta }$, where $B_{0\alpha \beta }$ is defined in \eqref{B0-ab};} \label{2 of Lemma 4.12}
		\item {For any $0 < R < + \infty$,
			\begin{equation*}
				\sup \limits_{n \geq 1} (1+|\xi|^2 )^{-1} \int_{|z-\xi |\leq R } A_{0\alpha \beta ,n } (z) dz \rightarrow 0 ,
		\end{equation*}	
where $A_{0 \alpha \beta , n} (z) = \int_{ \mathbb{S}^2 } B_{0 \alpha \beta, n} (z, \omega) \d \omega$. }\label{3 of Lemma 4.12}
	\end{enumerate}
\end{lemma}	

\begin{equation}\label{Q L}
	\begin{aligned}
		\tilde{Q}^\pm _{\alpha ,n } (f,f) &=
		\begin{cases}
			( N_n (f) )^{-1} Q^\pm _{\alpha ,n } (f,f)(\xi ), \quad &\alpha \in \{1,\cdots, s_0 \},\\
			( N_n (f) )^{-1} Q^\pm _{\alpha ,n } (f,f)(\xi ,I ), \quad &\alpha \in \{s_0+1,\cdots, s \},
		\end{cases} \\
			\tilde{Q}_{\alpha ,n } (f,f) &=
		\begin{cases}
			( N_n (f) )^{-1} Q_{\alpha ,n } (f,f)(\xi ), \quad &\alpha \in \{1,\cdots, s_0 \},\\
			( N_n (f) )^{-1} Q_{\alpha ,n } (f,f)(\xi ,I ), \quad &\alpha \in \{s_0+1,\cdots, s \},
		\end{cases}\\
		\tilde{L} _{\alpha ,n } (f) &=
		\begin{cases}
			( N_n (f) )^{-1} L _{\alpha ,n } (f)(\xi ), \quad & \ \ \ \alpha \in \{1,\cdots, s_0 \},\\
			( N_n (f) )^{-1} L _{\alpha ,n } (f)(\xi ,I ), \quad & \ \ \ \alpha \in \{s_0+1,\cdots, s \} \,, 
		\end{cases}
	\end{aligned}
\end{equation}

\begin{equation}\label{Nnf}
  \begin{aligned}
    N_n (f) = 1 + \tfrac{1}{n} \sum_{\alpha = 1}^{s_0} \int_{\R^3} | f_\alpha | d \xi + \tfrac{1}{n} \sum_{\alpha = s_0 + 1}^{s} \iint_{\R^3\times \R_+} | f_\alpha | d\xi dI \,. 
  \end{aligned}
\end{equation}

\begin{lemma}[Existence, Uniqueness and Positivity]\label{Lemma 4.14}
Consider the approximated initial data $f_{\alpha, 0}^n$ constructed in Lemmas \ref{Lemma 4.8}-\ref{Lemma 4.9} and the approximated collision operator $\tilde{Q}_{n} = ( \tilde{Q}_{\alpha, n} )_{\alpha \in \{ 1, \cdots, s \}}$ constructed in \eqref{Q L}. Then the approximated problem \eqref{approximation} admits a unique distributional solution $f^n = ( f^n_\alpha )_{\alpha \in \{ 1, \cdots, s \}}$ satisfying that for any fixed $T>0$,
\begin{enumerate}
\item{for $\alpha \in \{1,\cdots,s_0\}, f^n_\alpha \in C([0,+\infty ); L^1(\R^3 \times \R^3 ))$ and
\begin{equation*}
	f^n_\alpha (t,x,\xi ) \geq \tfrac{1}{n} \exp [ -c_0 (|x|^2+|\xi |^2 ) ] \quad \textrm{for all } \ (t,x,\xi ) \in [0, T] \times \R^3 \times \R^3 ;
	\end{equation*}}
\item{for $\alpha \in \{s_0+1,\cdots,s\}, f^n_\alpha \in C([0,+\infty ); L^1(\R^3 \times \R^3 \times \R_+))$,
	\begin{equation*}
		f^n_\alpha (t,x,\xi , I ) \geq \tfrac{1}{n} \exp [ -c_0 (|x|^2+|\xi |^2 + I ) ] \quad \textrm{for all } \ (t,x,\xi , I ) \in [0, T] \times \R^3 \times \R^3 \times \R_+ \,.
\end{equation*}}	
\end{enumerate}
Here the constant $c_0 = c_0 (n,T)>0 $.
\end{lemma}

\begin{lemma}\label{Lemma 4.17}
Assume that $f^n$ is the distributional solution to the approximated problem \eqref{approximation} constructed in Lemma \ref{Lemma 4.14}. Then it further enjoys the following properties:
\begin{enumerate}
	\item {Mass, momentum, energy and angular momentum are all conserved. Hence, for all $t \geq 0$ and $\alpha \in \{1,\cdots,s_0\}$, we have
		\begin{equation*}
			\iint_{\R^3 \times \R^3 } f^n_\alpha 
			\left (\begin{array}{c}
				1 \\[2mm]
				\xi \\ [2mm]
				|\xi|^2 \\ [2mm]
				|x-t\xi|^2 \\[2mm]
			\end{array}\right ) dxd\xi 
			= \iint_{\R^3 \times \R^3 } f_{\alpha ,0}^n 
			\left (\begin{array}{c}
				1 \\[2mm]
				\xi \\ [2mm]
				|\xi|^2 \\ [2mm]
				|x|^2 \\[2mm]
			\end{array}\right )dxd\xi ;
	\end{equation*}
For all $t>0$ and $\alpha \in \{s_0+1,\cdots,s\}$, we have
\begin{equation*}
	\iiint_{\R^3 \times \R^3 \times \R_+ } f^n_\alpha 
	\left (\begin{array}{c}
		1 \\[2mm]
		\xi \\ [2mm]
		\frac{1}{2}m_\alpha|\xi|^2+I \\ [2mm]
		|x-t\xi|^2 \\[2mm]
	\end{array}\right ) dxd\xi dI 
	= \iiint_{\R^3 \times \R^3 \times \R_+ } f_{\alpha ,0}^n 
	\left (\begin{array}{c}
		1 \\[2mm]
		\xi \\ [2mm]
		\frac{1}{2}m_\alpha|\xi|^2+I\\ [2mm]
		|x|^2 \\[2mm]
	\end{array}\right )dxd\xi dI .
	\end{equation*}}\label{1 of Lemma 4.17}
	\item{Entropy identity is true. Hence for all $t \geq 0$,
		\begin{equation}\label{entro}
			\begin{aligned}
				H(f^n_{0}) & = H(f^n)(t)+ \tfrac{1}{4}\sum \limits_{\alpha =1}^{s_0} \int_{0}^{t} d \tau \iint_{ \R^3 \times \R^3 } \tilde{e}_\alpha^n(
\tau,x,\xi ) dxd\xi\\
				& +\tfrac{1}{4}\sum \limits_{\alpha =s_0+1}^{s}\int_{0}^{t} d \tau \iiint_{ \R^3 \times \R^3 \times \R_+} \tilde{e}_\alpha^n(\tau,x,\xi ,I) dxd\xi dI ,
			\end{aligned}
		\end{equation}
where the density of entropy dissipation
\begin{equation}\label{e na}
	\begin{aligned}
	\tilde{e}_\alpha^n &= (N_n (f^n))^{-1} \sum \limits_{\beta =1}^s \int_{(\R^3\times \R_+)^3} (\tfrac{I^{\delta(\alpha)/2-1} I_* ^{\delta(\beta)/2-1}{f_\alpha^n}' {f_{\beta_*}^n}'  } {f^n_\alpha f^n_{\beta_* } (I')^{\delta(\alpha)/2-1} (I'_* )^{\delta(\beta)/2-1}} - 1)\\
&\times \log (\tfrac{I^{\delta(\alpha)/2-1} I_* ^{\delta(\beta)/2-1} {f_\alpha^n}' {f_{\beta_*}^n}' }{f^n_\alpha f^n_{\beta_* } (I')^{\delta(\alpha)/2-1} (I'_* )^{\delta(\beta)/2-1} } ) \tfrac{f^n_\alpha f^n_{\beta_* } }{I^{\delta(\alpha)/2-1} I_* ^{\delta(\beta)/2-1} }W_{\alpha\beta,n }d\xi_* d\xi'd\xi'_* dI_* dI' dI'_*,
	\end{aligned}
\end{equation}
the $H$-functional $H(f)$ is defined in \eqref{Hf}, and the transition probabilities $W_{\alpha \beta ,n}$ corresponds to the collision kernels $B_{2 \alpha \beta ,n}$, $B_{1 \alpha \beta, n}$, $B_{1 \beta \alpha, n}$ and $B_{0 \alpha \beta, n}$ constructed in Lemmas \ref{Lemma 4.10}-\ref{Lemma 4.11}-\ref{Lemma 4.12}. }\label{2 of Lemma 4.17}

\item{For all $\alpha \in \{1,\cdots,s_0\}$ and $t \geq 0$, we have
	\begin{equation}\label{FB2}{\small
		\begin{aligned}
  \iint_{ \R^3 \times \R^3} f^n_\alpha (t,x,\xi)(1+|x|^2+|\xi|^2)dxd\xi \leq & \iint_{ \R^3 \times \R^3} f^n_{\alpha,0}(x,\xi)(1+2|x|^2+(2t^2+1)|\xi|^2)dxd\xi,
  \end{aligned}}
	\end{equation}
and
 \begin{equation}\label{FWC2}
		\begin{aligned}
			\sup \limits_{t\geq 0}&\iiint_{ \R^3 \times \R^3} f^n_\alpha(t,x,\xi)|\log f^n_\alpha| dxd\xi +\tfrac{1}{4}\int_{0}^{+\infty} dt \iint_{ \R^3 \times \R^3} \tilde{e}^n_\alpha dxd\xi \\
		   \leq & 2 \iiint_{ \R^3 \times \R^3} f^n_{\alpha,0}(x,\xi) (|\log f^n_{\alpha,0}|+ |x|^2+ |\xi|^2 )dxd\xi +C_1;
		\end{aligned}
	\end{equation}
	For all $\alpha \in \{s_0+1,\cdots,s\}$ and $t \geq 0$, we have
	\begin{equation}\label{FB3}
		\begin{aligned}
   &\iint_{ \R^3 \times \R^3\times \R_+ } f^n_\alpha (t,x,\xi,I)(1+|x|^2+ \tfrac{1}{2} m_\alpha |\xi|^2+I)dxd\xi dI\\
  \leq & \iint_{ \R^3 \times \R^3\times \R_+ } f^n_{\alpha,0}(x,\xi,I ) \{ 1 + 2 |x|^2 + ( 2 t^2 + \tfrac{1}{2} m_\alpha ) |\xi|^2 + (\tfrac{4}{m_\alpha} t^2 + 1) I )\} dxd\xi dI,
\end{aligned}
\end{equation}
and
 \begin{equation}\label{FWC3}{\small
		\begin{aligned}
			\sup \limits_{t\geq 0}& \iiint_{ \R^3 \times \R^3\times \R_+} f^n_\alpha(t,x,\xi,I) |\log (I^{1-\delta(\alpha)/2}f^n_\alpha)| dxd\xi dI +\tfrac{1}{4}\int_{0}^{+\infty} dt \iiint_{ \R^3 \times \R^3\times \R_+} \tilde{e}^n_\alpha dxd\xi dI\\
		\leq & 2 \iiint_{ \R^3 \times \R^3\times \R_+} f^n_{\alpha,0}(x,\xi,I )\{|\log (I^{1-\delta(\alpha)/2}f^n_{\alpha,0})|+ M_\alpha(|x|^2+m_\alpha |\xi|^2 +2I)\}dxd\xi dI+C_1,
		\end{aligned}}
\end{equation}
where $M_\alpha = \max \{1,\frac{1}{m_\alpha}\}$ and the constant $C_1>0$ is independent of $n$.}\label{3 of Lemma 4.17}
\end{enumerate}
\end{lemma}

\begin{equation}\label{Trans-fab}
	\begin{aligned}	
		& \iiint_{\R^3 \times B_R \times (0,R)} L_{ \alpha \beta ,n} (f^n) (t,x,\xi, I) dxd\xi dI \\
= & \iiint_{\R^3 \times \R^3 \times \R_+}f^n_{\beta} (t,x,\xi_*, I_*) ( \iint_{B_R \times (0,R)}  A_{2\alpha \beta ,n}(\xi-\xi_*,I+I_*) d\xi dI ) dxd\xi_* dI_* \\
			\leq & \epsilon \iiint_{\R^3 \times \R^3 \times \R_+} f^n_{\beta} (t,x ,\xi_*, I_*) (1+|\xi_*|^2 +I_*) dxd\xi_* dI_* + C_\epsilon \iiint_{\R^3 \times \R^3 \times \R_+} f^n_{\beta} (t,x, \xi_*, I_*) dxd\xi_* dI_*.
		\end{aligned}
	\end{equation}

\begin{lemma}[Averaged velocity-internal energy lemma]\label{theo-6.2}
Let $g^n (t,x, \xi, I)$ weakly converge to $g (t,x, \xi, I)$ in $ L^1_{loc} ((0,T)\times \R^3 \times \R^3 \times \R_+) $, and $f^n (t,x, \xi, I)$ weakly converge to $f (t, x, \xi, I)$ in $L^1 ( (0,T)\times \R^3 \times \R^3 \times \R_+ )$. Moreover, they satisfy
	\begin{equation*}
		\tfrac{\partial f^n}{\partial t} +\xi \cdot \nabla_x f^n = g^n \quad \text{in}\ \mathscr{D}' ((0,T) \times \R^3 \times \R^3 \times \R_+) .
	\end{equation*}
Set $\{\varphi^n (t, x, \xi, I) \}_{n\geq 1} \subset L^\infty ((0,T)\times \R^3 \times \R^3 \times \R_+)$ converge to $\varphi (t, x, \xi, I)$ $a.e.$ $(0,T)\times \R^3 \times \R^3 \times \R_+$. Then the sequence $\iint_{\R^3\times \R_+} f^n\varphi^n d\xi dI$ strongly converges to $\iint_{\R^3\times \R_+} f \varphi d\xi dI$ in $L^1((0,T)\times \R^3)$.
\end{lemma}

\begin{lemma}[Averaged velocity lemma]\label{theo-6.2.1}
Let $g^n (t,x, \xi)$ weakly converge to $g (t, x, \xi)$ in $L^1_{loc} ( (0,T) \times \R^3 \times \R^3 )$, and $ f^n (t, x, \xi) $ weakly converge to $f (t, x, \xi)$ in $L^1 ( (0,T) \times \R^3 \times \R^3 )$. Moreover, they satisfy
	\begin{equation*}
		\tfrac{\partial f^n}{\partial t} +\xi \cdot \nabla_x f^n = g^n \quad \text{in}\ \mathscr{D}' ((0,T) \times \R^3 \times \R^3 ) .
	\end{equation*}
Set $\{\phi^n (t,x, \xi)\}_{n\geq 1} \subset L^\infty ((0,T)\times \R^3 \times \R^3 )$ converge to $ \phi (t, x, \xi) $ $a. e.$ $(0,T) \times \R^3 \times \R^3$. Then the sequence $\int_{\R^3} f^n\phi^n d\xi $ strongly converges to $ \int_{\R^3 } f \phi d\xi $ in $L^1((0,T)\times \R^3)$.
\end{lemma}

\begin{lemma}\label{theo-6.4}

Let the bounded sequence $\{ \varphi^n (t, x, \xi, I, e)\}_{n\geq 1}$ in $L^\infty ((0,T)\times \R^3 \times \R^3\times \R_+ ; L^1(E))$ satisfy that there exists $\varphi (t, x, \xi, I, e) \in L^\infty ((0,T)\times \R^3 \times \R^3 \times \R_+; L^1(E))$ such that 
\begin{equation*}
\lim \limits_{n\rightarrow \infty} \Vert \varphi^n -\varphi \Vert_{L^1(E)}(t,x,\xi,I)=0 \ a.e. \ (0,T)\times \R^3 \times \R^3\times \R_+.
\end{equation*}
Besides, $f^n (t, x, \xi, I)$ and $g^n (t, x, \xi, I)$ are given in Lemma \ref{theo-6.2}. Then the sequence $\iint_{\R^3\times \R_+} f^n\varphi^n d\xi dI$ strongly converges to $ \iint_{\R^3\times \R_+} f\varphi d\xi dI $ in $L^1((0,T)\times \R^3\times E)$.
\end{lemma}

\begin{lemma}\label{theo-6.4.1}
Let the bounded sequence $\{	\phi^n (t, x, \xi, e)\}_{n\geq 1}$ in $L^\infty ((0,T)\times \R^3 \times \R^3 ; L^1(E))$ satisfy that there exists $\phi \in L^\infty ((0,T)\times \R^3 \times \R^3 ; L^1(E))$ such that 
\begin{equation*}
	\lim \limits_{n\rightarrow \infty} \Vert \phi^n -\phi \Vert_{L^1(E)}(t,x,\xi)=0 \ a.e. \ (0,T)\times \R^3 \times \R^3.
\end{equation*}
Besides, $f^n (t, x, \xi)$ and $g^n (t, x, \xi)$ are given in Lemma \ref{theo-6.2.1}. Then the sequence $\int_{\R^3} f^n\phi^n d\xi $ strongly converges to $ \int_{\R^3 } f \phi d\xi $ in $L^1((0,T)\times \R^3\times E)$.
\end{lemma}

\begin{lemma}\label{converge f,g}
Let $T >0$ be fixed. 
\begin{enumerate}
\item {For $\alpha \in \{1,\cdots,s_0\}$, one has
	\begin{equation}\label{eq6.77}
	\begin{aligned}
	\lim \limits_{\delta \rightarrow 0^+} \sup \limits_{n\geq 1 } \sup \limits_{t\in (0,T)} \Vert f^n_\alpha -g^n_{\delta \alpha} \Vert_{L^1(\R^3 \times \R^3 )}=0, \ \lim \limits_{\delta \rightarrow 0^+} \Vert f_\alpha - \tilde{g}_{\delta \alpha} \Vert_{L^1( (0, T) \times \R^3 \times \R^3  )}=0;
	\end{aligned}
	\end{equation}}
\item{For $\alpha \in \{s_0+1,\cdots,s\}$, one has
\begin{equation}\label{eq6.76}
		\begin{aligned}
		\lim \limits_{\delta \rightarrow 0^+} \sup \limits_{n\geq 1 } \sup \limits_{t\in (0,T)} \Vert f^n_\alpha -g^n_{\delta \alpha} \Vert_{L^1(\R^3 \times \R^3\times \R_+ )}=0, \ \lim \limits_{\delta \rightarrow 0^+} \Vert f_\alpha - \tilde{g}_{\delta \alpha} \Vert_{L^1( (0, T) \times \R^3 \times \R^3\times \R_+)}=0.
	\end{aligned}
\end{equation}}
\end{enumerate}
\end{lemma}

\begin{lemma}\label{converge L}
	Let $T,R>0$ fixed. 
	\begin{enumerate}
		\item {For $\alpha \in \{1,\cdots,s_0\}$, one has
			\begin{equation}
				\lim \limits_{n\rightarrow \infty} \tilde{L}_{\alpha,n} (f^n) = L_\alpha(f) \quad \text{strongly in } L^1 ((0,T)\times \R^3\times B_R);
		\end{equation}}
		\item{For $\alpha \in \{s_0+1,\cdots,s\}$, one has
			\begin{equation}
			\lim \limits_{n\rightarrow \infty} \tilde{L}_{\alpha,n} (f^n) = L_\alpha(f) \quad \text{strongly in } L^1 ((0,T)\times \R^3\times B_R\times (0,R)).
		\end{equation}}
	\end{enumerate}
\end{lemma}

\begin{proof}

Let $E=B_R,\ F=B_R\times (0,R)$ and $r > 0$ fixed. Consider the variables $t \in (0, T)$, $x \in \R^3$, $\xi \in B_R$, $I \in (0, R)$, $\xi_* \in \R^3$ and $I_* \in \R_+$. Denote by
\begin{equation*}
	\begin{aligned}
		\varphi^n_1 (t, x, \xi_*, \xi) & = A_{0\alpha \beta ,n}(\xi -\xi_* ) \mathbf{1}_{|\xi_* |\leq r} \in L^\infty ((0,T)\times \R^3 \times \R^3; L^1(E)), \\
		\varphi^n_2 (t, x, \xi_*, I_*,\xi) & = A_{1\alpha \beta ,n}(\xi -\xi_* ,I_* ) \mathbf{1}_{|\xi_* |\leq r} \mathbf{1}_{0 \leq I_* \leq r} \in L^\infty ((0,T)\times \R^3 \times \R^3 \times \R_+; L^1(E)),\\
		\varphi^n_3 (t, x, \xi_*, \xi, I) & = A_{1\beta \alpha ,n}(\xi -\xi_* ,I) \mathbf{1}_{|\xi_* |\leq r} \in L^\infty ((0,T)\times \R^3 \times \R^3; L^1(F)), \\
		\varphi^n_4 (t, x, \xi_*, I_*, \xi, I) & = A_{2\alpha \beta ,n}(\xi -\xi_* ,I+I_* ) \mathbf{1}_{|\xi_* |\leq r} \mathbf{1}_{0 \leq I_* \leq r} \in L^\infty ((0,T)\times \R^3 \times \R^3\times \R_+; L^1(F)) \,.
	\end{aligned}
\end{equation*}
Similarly, let
\begin{equation*}
	\begin{aligned}
		& \varphi_1 (t, x, \xi_*, \xi) = A_{0\alpha \beta }(\xi -\xi_* ) \mathbf{1}_{|\xi_* |\leq r}, \quad \quad \ \ \, \varphi_2 (t, x, \xi_*, I_*, \xi ) = A_{1\alpha \beta }(\xi -\xi_* ,I_* ) \mathbf{1}_{|\xi_* |\leq r} \mathbf{1}_{0 \leq I_* \leq r}, \\
		&\varphi_3 (t, x, \xi_*, \xi, I) = A_{1\beta \alpha}(\xi -\xi_* ,I) \mathbf{1}_{|\xi_* |\leq r}, \quad \varphi_4 (t, x, \xi_*, I_*, \xi, I ) = A_{2\alpha \beta }(\xi -\xi_* ,I+I_* ) \mathbf{1}_{|\xi_* |\leq r} \mathbf{1}_{0 \leq I_* \leq r}.
	\end{aligned}
\end{equation*}
Together with Lemmas \ref{Lemma 4.10}-\ref{Lemma 4.11}-\ref{Lemma 4.12}, one obtains  
\begin{equation*}
	\begin{aligned}
		& \lim \limits_{n\rightarrow \infty } \Vert \varphi^n_1 -\varphi_1 \Vert_{L^1(B_R)}=0, \quad a.e. \ (t,x,\xi_*) \in (0,T) \times \R^3 \times \R^3, \\
		& \lim \limits_{n\rightarrow \infty } \Vert \varphi^n_2 -\varphi_2 \Vert_{L^1(B_R)}=0, \quad a.e. \ (t,x,\xi_* ,I_*) \in (0,T)\times \R^3 \times \R^3 \times \R_+, \\
		& \lim \limits_{n\rightarrow \infty } \Vert \varphi^n_3 -\varphi_3 \Vert_{L^1(B_R\times (0,R))}=0, \quad a.e. \ (t,x,\xi_*) \in (0,T)\times \R^3 \times \R^3, \\
		& \lim \limits_{n\rightarrow \infty } \Vert \varphi^n_4 -\varphi_4 \Vert_{L^1(B_R\times (0,R) )}=0, \quad a.e. \ (t,x,\xi_* ,I_*) \in (0,T)\times \R^3 \times \R^3 \times \R_+ . 
	\end{aligned}
\end{equation*}
Then Lemma \ref{theo-6.4.1} shows that
\begin{equation*}{\small
	\begin{aligned}
		\lim \limits_{n\rightarrow \infty } \int_{|\xi_* |\leq r} g^n_{\delta \beta_* } A_{0\alpha \beta ,n}(\xi -\xi_* ) d\xi_* &=\int_{|\xi_* |\leq r} \tilde{g}_{\delta \beta_*} A_{0\alpha \beta } (\xi -\xi_* )d\xi_*, \\
		\lim \limits_{n\rightarrow \infty } \int_{|\xi_* |\leq r} \int_{0\leq I_* \leq r} g^n_{\delta \beta_*} A_{1\alpha \beta ,n}(\xi -\xi_* ,I_* ) d\xi_* dI_*&=\int_{|\xi_* |\leq r} \int_{0\leq I_* \leq r} \tilde{g}_{\delta \beta_*} A_{1\alpha \beta } (\xi -\xi_* ,I_* )d\xi_* dI_*
	\end{aligned}}
\end{equation*}
strongly in $L^1((0,T)\times \R^3 \times B_R )$. Moreover, Lemma \ref{theo-6.4} indicates that
\begin{equation*}
	\begin{aligned}
		\lim \limits_{n\rightarrow \infty } &\int_{|\xi_* |\leq r} g^n_{\delta \beta_*} A_{1\beta \alpha ,n}(\xi -\xi_* ,I) d\xi_* =\int_{|\xi_* |\leq r} \tilde{g}_{\delta \beta_*} A_{1\beta \alpha }(\xi -\xi_* ,I) d\xi_*, \\
		\lim \limits_{n\rightarrow \infty } &\int_{|\xi_* |\leq r} \int_{0\leq I_* \leq r} g^n_{\delta \beta_*} A_{2\alpha \beta ,n}(\xi -\xi_* ,I+I_* ) d\xi_* dI_*\\
		=&\int_{|\xi_* |\leq r} \int_{0\leq I_* \leq r} \tilde{g}_{\delta \beta_*} A_{2\alpha \beta }(\xi -\xi_* ,I+I_* ) d\xi_* dI_*
	\end{aligned}
\end{equation*}
strongly in $L^1((0,T)\times \R^3 \times B_R \times (0,R))$.

Recall that $0 < g^n_{\delta \alpha} < f^n_\alpha$ for $\alpha \in \{ 1, \cdots, s \}$. For the cases $\alpha, \beta \in \{s_0+1,\cdots ,s\}$, following the similar arguments in \eqref{Trans-fab}, one knows that for all $\epsilon>0$ there exists $C_\epsilon >0$ such that
\begin{equation*}
	\begin{aligned}
		&\Vert \int_{|\xi_* |\leq r} \int_{0\leq I_* \leq r} (f^n_{ \beta_*} -g^n_{\delta \beta_*}) A_{2\alpha \beta ,n}(\xi -\xi_* ,I+I_* ) d\xi_* dI_* \Vert_{L^1((0,T)\times \R^3 \times B_R \times (0,R))} \\
		\leq & \epsilon \int_{0}^{T} dt \iiint_{\R^3 \times \R^3 \times \R_+} (f^n_{\beta *} -g^n_{\delta \beta * } ) (1+|\xi_*|^2+I_*) dx d\xi_* dI_* \\
		+&C_\epsilon  \int_{0}^{T} dt \iiint_{\R^3 \times \R^3 \times \R_+}| f^n_{\beta *} -g^n_{\delta \beta * } |dx d\xi_* dI_* .
	\end{aligned}
\end{equation*}
By Lemma \ref{converge f,g}, one knows that as $\delta \to 0^+$,
\begin{equation*}{\small
  \begin{aligned}
    C_\epsilon  \int_{0}^{T} dt \iiint_{\R^3 \times \R^3 \times \R_+}| f^n_{\beta *} -g^n_{\delta \beta * } |dx d\xi_* dI_* \leq C_\epsilon T \sup_{n \geq 1} \sup_{ t \in (0, T) } \iiint_{\R^3 \times \R^3 \times \R_+}| f^n_{\beta *} -g^n_{\delta \beta * } |dx d\xi_* dI_* \to 0 \,.
  \end{aligned}}
\end{equation*}
Moreover, by Lemma \ref{Lemma 4.17},
\begin{equation*}
  \begin{aligned}
    & \epsilon \int_{0}^{T} dt \iiint_{\R^3 \times \R^3 \times \R_+} (f^n_{\beta *} -g^n_{\delta \beta * } ) (1+|\xi_*|^2+I_*) dx d\xi_* dI_* \\
    \leq & 2 \epsilon T \sup_{ t \in [0, T] } \iiint_{\R^3 \times \R^3 \times \R_+} f^n_{\beta *} (1+|\xi_*|^2+I_*) dx d\xi_* dI_* \leq C_T \epsilon \to 0 \ (\epsilon \to 0) \,.
  \end{aligned}
\end{equation*}
Consequently, as $\delta \to 0^+$,
\begin{equation*}
  \begin{aligned}
    \Vert \int_{|\xi_* |\leq r} \int_{0\leq I_* \leq r} (f^n_{ \beta_*} -g^n_{\delta \beta_*}) A_{2\alpha \beta ,n}(\xi -\xi_* ,I+I_* ) d\xi_* dI_* \Vert_{L^1((0,T)\times \R^3 \times B_R \times (0,R))} \to 0 \,.
  \end{aligned}
\end{equation*}
Similarly, as $\delta \to 0^+$,
\begin{equation*}
	\Vert \int_{|\xi_* |\leq r} \int_{0\leq I_* \leq r} (f_{ \beta_*} -\tilde{g}_{\delta \beta_*}) A_{2\alpha \beta ,n}(\xi -\xi_* ,I+I_* ) d\xi_* dI_* \Vert_{L^1((0,T)\times \R^3 \times B_R \times (0,R))} \rightarrow  0.
\end{equation*}
Therefore, for the cases $ \alpha, \beta \in \{ s_0 + 1, \cdots, s \} $, as $n \rightarrow + \infty $ and $ \delta \rightarrow 0^+$,
	\begin{align}\label{K1}
		\no &\Vert \int_{|\xi_* |\leq r} \int_{0\leq I_* \leq r} (f^n_{ \beta_*}-f_{\beta_*}) A_{2\alpha \beta ,n}(\xi -\xi_* ,I+I_* ) d\xi_* dI_*  \Vert_{L^1((0,T)\times \R^3 \times B_R \times (0,R))} \\
		\no \leq & \Vert \int_{|\xi_* |\leq r} \int_{0\leq I_* \leq r} (g^n_{\delta \beta_*}-\tilde{g}_{\delta \beta_*}) A_{2\alpha \beta ,n} (\xi -\xi_* ,I+I_* )d\xi_* dI_* \Vert_{L^1((0,T)\times \R^3 \times B_R \times (0,R))} \\
		\no + & \sup_{n \geq 1} \Vert \int_{|\xi_* |\leq r} \int_{0\leq I_* \leq r} (f^n_{ \beta_*} -g^n_{\delta \beta_*}) A_{2\alpha \beta ,n} (\xi -\xi_* ,I+I_* )d\xi_* dI_* \Vert_{L^1((0,T)\times \R^3 \times B_R \times (0,R))} \\
		+ &\Vert \int_{|\xi_* |\leq r} \int_{0\leq I_* \leq r} (f_{\beta_*} -\tilde{g}_{\delta \beta_*}) A_{2\alpha \beta,n }(\xi -\xi_* ,I+I_* ) d\xi_* dI_* \Vert_{L^1((0,T)\times \R^3 \times B_R \times (0,R))} \rightarrow 0.
	\end{align}

On the other hand, following the similar arguments in \eqref{Trans-fab} and combining with the uniform bounds in Lemma \ref{Lemma 4.17}, one knows that for any $\epsilon > 0$ there is a $C_\epsilon > 0$ such that
\begin{equation}\label{K2}
	\begin{aligned}
		&\Vert \iint_{ \R^3 \times \R_+} f^n_{ \beta *} A_{2\alpha \beta ,n}(\xi -\xi_* ,I+I_* ) ( \mathbf{1}_{|\xi_*|>r}+ \mathbf{1}_{I_*>r})d\xi_* dI_* \Vert_{L^1((0,T)\times \R^3 \times B_R \times (0,R))} \\
		\leq & \tfrac{C_\epsilon}{r^2} \int_0^T dt \iiint_{ \R^3 \times \R^3 \times \R_+} f^n_{\beta *} |\xi_*|^2 dxd\xi_* dI_*  + \tfrac{C_\epsilon}{r} \int_0^T dt \iiint_{ \R^3 \times \R^3 \times \R_+} f^n_{\beta *} \cdot I_* dxd\xi_* dI_* \\
		+ & \epsilon  \int_0^T dt \iiint_{ \R^3 \times \R^3 \times \R_+} f^n_{\beta *} (1+|\xi_*|^2+I_*) d\xi_* dI_* \rightarrow 0 \quad (r \rightarrow \infty, \epsilon \rightarrow 0^+).
	\end{aligned}
\end{equation}
Similarly, for $\alpha, \beta \in \{ s_0 + 1, \cdots, s \}$,
\begin{equation}\label{K33}
	\lim_{r \to + \infty} \Vert \iint_{ \R^3 \times \R_+} f_{ \beta *} (A_{2\alpha \beta }+ A_{2\alpha \beta,n }) ( \mathbf{1}_{|\xi_*|>r}+ \mathbf{1}_{I_*>r})d\xi_* dI_* \Vert_{L^1((0,T)\times \R^3 \times B_R \times (0,R))} = 0.
\end{equation}
Therefore, the convergence \eqref{K1}-\eqref{K2}-\eqref{K33} imply that as $n \to + \infty$ and $r \to + \infty$,
\begin{equation}\label{K4}
	\begin{aligned}
		&\Vert L_{ \alpha \beta ,n} (f^n ) -L_{ \alpha \beta} (f) \Vert_{L^1((0,T)\times \R^3 \times B_R \times (0,R))}\\
		\leq &\Vert \int_{|\xi_* |\leq r} \int_{0\leq I_* \leq r} (f^n_{ \beta_*}-f_{\beta_*}) A_{2\alpha \beta ,n}(\xi -\xi_* ,I+I_* ) d\xi_* dI_*  \Vert_{L^1((0,T)\times \R^3 \times B_R \times (0,R))} \\
		+&\Vert \iint_{ \R^3 \times \R_+} f^n_{ \beta_*} A_{2\alpha \beta ,n} ( \mathbf{1}_{|I_*|>r}+ \mathbf{1}_{I_*>r})d\xi_* dI_* \Vert_{L^1((0,T)\times \R^3 \times B_R \times (0,R))} \\
		+& \Vert \iint_{ \R^3 \times \R_+} f_{ \beta_*} (A_{2\alpha \beta }+ A_{2\alpha \beta,n }) ( \mathbf{1}_{|I_*|>r}+ \mathbf{1}_{I_*>r})d\xi_* dI_* \Vert_{L^1((0,T)\times \R^3 \times B_R \times (0,R))} \rightarrow 0 
	\end{aligned}
\end{equation}
for $\alpha, \beta \in \{ s_0 + 1, \cdots, s \}$.

Similarly, for $ \alpha \in \{ s_0 + 1, \cdots, s \} $, $\beta \in \{ 1, \cdots, s_0 \}$, 
\begin{equation}\label{K5}
  \begin{aligned}
    \lim_{n \to + \infty} L_{ \alpha \beta ,n} (f^n ) = L_{ \alpha \beta} (f) \quad \textrm{strongly in } \ L^1 ( (0, T) \times \R^3 \times B_R \times (0, R) ) \,.
  \end{aligned}
\end{equation}
Then, by \eqref{K4}-\eqref{K5}, we have
\begin{equation}\label{K6}
	\lim\limits_{n\rightarrow \infty } L_{\alpha ,n} (f^n) =L_{\alpha } (f) \quad \text{strongly in }\ L^1((0,T)\times \R^3 \times B_R\times (0,R))
\end{equation}
for all $\alpha \in \{s_0+1,\cdots,s\}$, which means that $ \sup_{n \geq 1} \| L_{\alpha ,n} (f^n) \|_{ L^1((0,T)\times \R^3 \times B_R\times (0,R)) } \leq C $. 

Recalling the definition of $N_n (f^n)$ in \eqref{Nnf}, we know that $N_n (f^n)^{-1}$ is bounded in $L^\infty ( (0,T) \times \R^3)$ and converges to $1$ almost everywhere. Together with \eqref{K6},
\begin{equation}\label{eq6.78}{\small
	\begin{aligned}
		& \Vert \tilde{L}_{\alpha ,n}(f^n) -L_\alpha (f) \Vert_{L^1((0,T)\times \R^3 \times B_R \times (0,R))} \\
		\leq & \Vert {L}_{\alpha ,n}(f^n) -L_\alpha (f) \Vert_{L^1((0,T)\times \R^3 \times B_R \times (0,R))} + \Vert \tilde{L}_{\alpha ,n}(f^n) - {L}_{\alpha ,n}(f^n) \Vert_{L^1((0,T)\times \R^3 \times B_R \times (0,R))}  .
	\end{aligned}}
\end{equation}


Then, by applying Egorov's theorem, for any $\epsilon > 0$, there exists a measurable subset $A \subset (0, T) \times B_R \times B_R \times (0, R)$ such that $ N_n (f^n)^{-1} $ converges uniformly to $1$ on $A$, and $\mu \left(((0, T) \times B_R \times B_R \times (0, R)) \setminus A\right) < \epsilon$. Thus,
\begin{equation}\label{Lan-La}
	\lim_{n\to\infty} \Vert \tilde{L}_{\alpha ,n}(f^n) - {L}_{\alpha ,n}(f^n) \Vert_{L^1((0,T)\times \R^3 \times B_R \times (0,R))}=0.
\end{equation}
Equation \eqref{K6} implies that $L_{\alpha,n}(f^n)$ is weakly compact in $L^1((0, T) \times \mathbb{R}^3 \times B_R \times (0, R))$. Then, as $\epsilon \to 0^+$ and $R \to \infty$,
\begin{equation}\label{K7}
	\|L_{\alpha,n}(f^n)\|_{L^1((0,T)\times B_R \times B_R \times (0,R)\setminus A)} + \|L_{\alpha,n}(f^n)\|_{L^1((0,T)\times (B_R)^c \times B_R \times (0,R))} \to 0.  
\end{equation}
Substituting \eqref{K6}, \eqref{Lan-La}, and \eqref{K7} into \eqref{eq6.78}, we obtain that the strong convergence holds.  
Similarly, for $\alpha \in \{1,\cdots,s_0\}$, we have 
\begin{equation*}
	\lim\limits_{n\rightarrow \infty } \tilde{L}_{\alpha ,n} (f^n) =L_{\alpha } (f) \quad \text{in}\ L^1((0,T)\times \R^3 \times B_R).
\end{equation*}
As a result, the proof of Lemma \ref{converge L} is finished.
\end{proof}

\end{document}
